\chapter{Formales}
\label{cha:form_DE}

In den folgenden Abschnitten wird auf verschiedene formale Aspekte eingegangen, die im Wesentlichen befolgt werden sollten. Die Vorgaben zur äußeren Form, der Gliederungstiefe und der Zitierweise sind unbedingt einzuhalten.

\section{Sprache}
\label{sec:sprache}

Wissenschaftliche Arbeiten am IWES können in deutscher oder englischer Sprache verfasst werden, wobei es für englischsprachige Arbeiten eine gesonderte Vorlage bzw. einen gesonderten Leitfaden gibt. Die Wahl der Sprache ist mit der Betreuerin bzw. dem Betreuer der Arbeit abzustimmen. In jedem Fall ist die gültige Rechtschreibung und Grammatik anzuwenden.

Der Sprachstil sollte einer wissenschaftlichen Ausarbeitung angemessen sein. Es ist Berichtssprache zu verwenden. Umgangssprachliche Formulierungen sind zu vermeiden. Wissenschaftliches Schreiben zeichnet sich durch einen sehr effizienten Sprachstil aus. Dieser ist jederzeit sachlich und so knapp wie möglich zu halten. Unnötige Füllwörter und Übergangsformulierungen sollten vermieden werden. Jedes Wort, das zum Verständnis der Aussage nicht unbedingt erforderlich ist, sollte weggelassen werden. Dabei ist natürlich trotzdem darauf zu achten, dass ein ansprechender Lesefluss gegeben ist.

Um die Fehlerdichte zu minimieren, sollten Tools für eine effektive Rechtschreib- und Grammatikprüfung vor der Abgabe der Arbeit Verwendung finden. Auch sprachliche Mängel können zu einer Abwertung der Ausarbeitung führen.

\section{Äußere Form}
\label{sec:form_DE}

Wissenschaftliche Ausarbeitungen sind einseitig, einspaltig und im Blocksatz bei 1,5-fachem Zeilenabstand zu erstellen. Als Schriftart ist eine serifenlose Schrift zu verwenden. Die Schriftgröße muss zwischen 11 pt und 12 pt gewählt werden. Alternativ kann, wenn die Arbeit mittels \LaTeX{} verfasst wird, die Standartschrift von \LaTeX{} bei 11 pt verwendet werden. Diese wird durch die Pakete \textbackslash{}usepackage[T1]\{fontenc\} und \textbackslash{}usepackage[utf8]\{inputenc\} auf den deutschen Zeichenvorrat ausgeweitet.

Abgesehen vom Deckblatt und der Titelseite, die eigene Formatierungen aufweisen, sind die folgenden Seitenränder einzuhalten:

\begin{itemize}
\item 2,5 cm (oben, links und rechts)
\item 2,0 cm (unten)
\end{itemize}

Kapitel müssen mit einer neuen Seite beginnen. Bei Unterkapiteln und Unter-Unterkapiteln wird der Seitenumbruch dagegen variabel gehalten.

\section{Gliederungsebenen, Kopfzeilen und Seitennummerierungen}
\label{sec:nummerierung}

Es dürfen nicht mehr als drei nummerierte Gliederungsebenen verwendet werden, also Kapitel, Unterkapitel und Unter-Unterkapitel. Sollte eine vierte Gliederungsebene erforderlich werden, ist die Ausarbeitung zu restrukturieren, um die Limitierung auf drei Gliederungsebenen einzuhalten.

Ab dem Titelblatt ist eine Kopfzeile vorzusehen. Am linken Rand der Kopfzeile werden die Kapitelnummer und die Kapitelbezeichnung aufgeführt. Die Seitennummerierung ist im rechten Teil der Kopfzeile anzuordnen. 

Die Nummerierung der einführenden Seiten (Eidesstattliche Erklärung, Aufgabenstellung, Kurzfassung und Inhaltsverzeichnis) erfolgt in römischen Zahlen (I, II, III, ...). Die römische Zählung beginnt mit der Titelseite (erste Papierseite). Ab den Verzeichnissen (Abbildungsverzeichnis etc.) ist die römische Nummerierung fortlaufend fortzuführen. 

\begin{addmargin}[1cm]{1cm}
	\textit{\textbf{Beispiel:} Endet das Inhaltsverzeichnis mit der römischen Seitenzahl \textit{VI}, ist die Seitennummerierung bei den Verzeichnissen mit der römischen Seitenzahl \textit{VII} weiterzuführen.}
\end{addmargin}

Der eigentliche Hauptteil des Textes, also alle Kapitel von der Einleitung bis zur Zusammenfassung und dem Ausblick, wird, mit 1 beginnend, mit arabischen Zahlen nummeriert (1, 2, 3, ...). Die Zählung beginnt mit der ersten Seite der Einleitung.

\section{Formeln, Abbildungen und Tabellen}

Die Anforderungen für die Nutzung und Einbindung von Formeln, Abbildungen und Tabellen werden in gesonderten Kapiteln dargelegt, s. die Abschnitte \ref{cha:Formeln}, \ref{cha:Images_DE} und \ref{cha:Tabellen}. Dort werden auch Beispiele für unterschiedliche Arten der Einbindung dargestellt. Gemeinsame Anforderungen werden im Folgenden definiert.

Formeln, Abbildungen und Tabellen sind zu nummerieren. Hierfür ist eine kapitelinterne, fortlaufende Nummerierung zu wählen. Das bedeutet, dass die jeweiligen Nummern mit der Kapitelnummer beginnen und nach einem folgenden Punkt die Nummerierung fortlaufend erfolgt. Die nach dem Punkt befindliche fortlaufende Nummerierung ist für Formeln, Abbildungen und Tabellen jeweils unabhängig voneinander und muss in jedem Kapitel wieder bei 1 beginnen. Das Prinzip ist in Tabelle \ref{tab:Nummerierung} zusammengefasst. Die Gliederungsebene innerhalb eines Kapitels wird bei der Nummerierung nicht berücksichtigt.

\begin{table}[!htbp]
	\small
	\centering
	\caption[Nummerierungsbeispiele für Formeln, Abbildungen und Tabellen]{Nummerierungsbeispiele für Formeln, Abbildungen und Tabellen}
	\begin{tabular}{P{2.5cm}P{1.5cm}P{1cm}P{2cm}P{1cm}P{1.5cm}P{1cm}}
		\toprule
		\textbf{Kapitel} & \textbf{Formel} & \textbf{Nr.} & \textbf{Abbildung} & \textbf{Nr.} & \textbf{Tabelle} & \textbf{Nr.} \\
		\midrule
		Kapitel 1 
		& 1. Formel & 1.1 & 1. Abbildung & 1.1 & 1. Tabelle & 1.1 \\
		& 2. Formel & 1.2 & 2. Abbildung & 1.2 & 2. Tabelle & 1.2 \\
		& ... & ... & ... & ... & ... & ... \\
		Kapitel 2 
		& 1. Formel & 2.1 & 1. Abbildung & 2.1 & 1. Tabelle & 2.1 \\
		& 2. Formel & 2.2 & 2. Abbildung & 2.2 & 2. Tabelle & 2.2 \\
		& ... & ... & ... & ... & ... & ... \\
		Kapitel 3
		& 1. Formel & 3.1 & 1. Abbildung & 3.1 & 1. Tabelle & 3.1 \\
		& 2. Formel & 3.2 & 2. Abbildung & 3.2 & 2. Tabelle & 3.2 \\
		... & ... & ... & ... & ... & ... & ... \\
		\bottomrule
	\end{tabular}%
	\label{tab:Nummerierung}%
\end{table}%



\section{Verweise und Fußnoten}
\label{sec:Verweise}

Innerhalb eines Kapitels sind oftmals Verweise auf andere Bestandteile der Arbeit notwendig. Verweise sollen mittels der jeweiligen Abschnittsnummer vorgenommen werden. Will man z.B. auf den Inhalt des Abschnittes "`Aufgabenstellung"' verweisen, so hat dies mit der Abschnittsnummer, in dem Fall \ref{sec:Aufgabe}, zu erfolgen. 

\begin{addmargin}[1cm]{1cm}
	\textit{\textbf{Beispiel:} Hinweise zum Einbinden der Aufgabenstellung finden sich in Abschnitt \ref{sec:Aufgabe}.}
\end{addmargin}

Derartige Verweise sind über die entsprechenden Befehle oder Funktionen der Textverarbeitungsprogramme zu erstellen. Ein händisches Eintragen der Nummern führt mit hoher Wahrscheinlichkeit zu Fehlern im Enddokument. 

Wird auf eine Seite verwiesen, die nicht Bestandteil eines Kapitels ist, kann alternativ auch durch die Angabe der Seitennummer auf dieses verwiesen werden. Dies ist aber im Allgemeinen zu vermeiden.

Wird auf Formeln, Abbildungen, Tabellen oder Anhänge Bezug genommen, so hat dies mittels der Angabe der Formel-, Abbildungs-, Tabellen- oder Anhangsnummer zu erfolgen (siehe hierzu auch Tabelle \ref{tab:Nummerierung}).

Fußnoten\footnote{Dies ist ein Beispiel einer Fußnote.} werden dann eingesetzt, wenn eine Aussage zwar für das Verständnis des Lesers wichtig, jedoch im Fließtext überflüssig oder fehlplatziert ist. Fußnoten werden über das gesamte Dokument hinweg durchgehend nummeriert. Die Fußzeile wird durch einen etwa 6,5 cm langen horizontalen Strich abgesetzt.

\section{Symbole und Abkürzungen}
\label{sec:Abk}

Symbole (auch Formelzeichen) und Abkürzungen sind innerhalb des Fließtextes dort einzuführen, wo sie zum ersten Mal verwendet werden. Für Symbole wird darauf in Abschnitt \ref{cha:Formeln} noch einmal eingegangen. Abkürzungen sind in Klammern hinter dem ausgeschriebenen Ausdruck einzuführen. 

\begin{addmargin}[1cm]{1cm}
	\textit{\textbf{Beispiel:} Dies ist der Leitfaden des Instituts für Windenergiesysteme (IWES) zum Anfertigen wissenschaftlicher Ausarbeitungen.}
\end{addmargin}



\section{Zitierrichtlinien}
\label{sec:Zitier}

Grundsätzlich sind alle aus anderen Veröffentlichungen wörtlich übernommenen Satzteile, Sätze oder Abschnitte sowie übernommene Abbildungen oder Tabellen zu zitieren (sogenannte direkte Zitate). Wörtlich übernommene Textbestandteile sind in Anführungsstrichen zu setzen. Das Zitat muss in dem Fall die Seiten- und Zeilenzahlen beinhalten. Sinngemäß übernommene Aussagen, Abbildungen, etc. sind ebenfalls durch die Angabe der Quelle als solche kenntlich zu machen (indirekte Zitate). Direkte Zitate sind im Ingenieurwesen äußerst unüblich und sollten vermieden werden. Indirekte Zitate sind grundsätzlich zu bevorzugen.

Eindeutige Quellenangaben sind beim wissenschaftlichen Arbeiten unumgänglich. Dadurch wird es der Leserin bzw. dem Leser ermöglicht, diese zu finden und somit die Herkunft der wissenschaftlichen Aussagen zu überprüfen, sofern es sich um nicht selbständig entwickeltes Gedankengut handelt.

Als Zitierstil ist in wissenschaftlichen Arbeiten am IWES eine numerische Zitierweise anzuwenden. Dafür ist die Nummer der jeweiligen Quelle im Literaturverzeichnis in eckigen Klammern hinter die zitierte Aussage zu setzen. 

\begin{addmargin}[1cm]{1cm}
	\textit{\textbf{Beispiel:} Eine detaillierte Beschreibung von Regelungskonzepten für große Windenergiesysteme findet sich in \cite{buch}.}
\end{addmargin}

Die Aufführung der Quellen im Literaturverzeichnis erfolgt in alphabetischer Reihenfolge. Es ist wichtig, dass die Formatierung von Einträgen im Literaturverzeichnis durchgängig konsistent ist. Dabei gibt es unterschiedliche Formatierungsvorschriften für unterschiedliche Quellenarten. Grundsätzlich sind überwiegend zitierfähige Quellen zu verwenden. Als zitierfähig gelten insbesondere Artikel in Fachzeitschriften \cite{artikel}, Fachbücher \cite{buch}, Beiträge in Konferenzbänden \cite{inp}, Teile eines Sammelbandes \cite{in}, technische Dokumentationen \cite{manual}, veröffentlichte Berichte einer Hochschule oder anderen Forschungsinstitution \cite{bericht} sowie Dissertationen \cite{thesis}. Nicht zitierfähige Quellen sind grundsätzlich zu vermeiden. Zwar kann insbesondere für Abbildungen nicht immer auf Internetquellen \cite{internet} verzichtet werden, dennoch sollten vorzugsweise Originalquellen (z. B. die Quellen von Wikipedia-Artikeln) verwendet werden. Die oben zitierten Referenzen stellen Beispiele dar. Die zugehörige Formatierung innerhalb des Literaturverzeichnisses sollte eingehalten werden.

\begin{addmargin}[1cm]{1cm}
	\textit{\textbf{Hinweis:} Vorlesungsunterlagen gelten als nicht zitierfähige Quellen, da sie als Allgemeinwissen behandelt werden. Im Zweifel sollte die in Vorlesungsunterlagen angegebene Originalliteratur zitiert werden. Auch Internetquellen sind als nicht zitierfähig einzustufen, da sie sich fortlaufend verändern und ggf. nicht überprüft werden können.}
\end{addmargin}


\section{Abgabe}
\label{sec:Abgabe}

Eine am IWES erstellte wissenschaftliche Arbeit ist in gedruckter und gebundener Form in einem weißen Umschlagkarton eingebunden abzugeben, auf dem das Deckblatt farbig aufgedruckt ist. 

\begin{addmargin}[1cm]{1cm}
	\textit{\textbf{Hinweis:} Gedruckte Fassungen in Ringbindung werden nicht angenommen, da eine Inventarisierung dann nicht möglich ist.}
\end{addmargin}


Die digitale Kopie der Arbeit ist als PDF-Datei per E-Mail an den Prüfer und die Betreuerin bzw. den Betreuer zu übersenden. Weitere elektronische Daten, einschließlich aller für die Reproduzierbarkeit relevanten Dateien (wie Quellen, Ergebnisse, Tabellen, Excel-Berechnungen, Simulationsergebnisse usw.), sollten über \textit{SeaFile} abgegeben werden. Der entsprechende Upload-Link wird auf Nachfrage von der Betreuerin bzw. dem Betreuer zur Verfügung gestellt.

Die Anzahl der abzugebenden gedruckten Fassungen ist der Aufgabenstellung zu entnehmen.



